\section{What remains}\label{sec:prog}

\subsection{Above logarithmic utility}

The one thing the programme set out to do and has not done is uniqueness at $\mu\in\{3,5\}$. The
argument of Section~\ref{sec:inc} does not fail there for want of propagation: the per-stage
multiplier $m_k=qA/(\Thorn+qA\kappa_kR)$ has a product below one over a life at every risk aversion we
tested, up to $\gamma=5$. What gives way is the certainty-equivalent side. The gradient sum $A$ is
one plus the variance of the risk tilt at $\gamma=1$ and only bounded by
$(1+\operatorname{Var})^{(\gamma+1)/(2\gamma)}$ above it, by Lyapunov, and both the loss in that
inequality and the variance itself grow with $\gamma$ --- the variance tenfold between $\gamma=1$ and
$\gamma=3$. The bound crosses the allowance between $\gamma=1$ and $\gamma=1.5$.

So the question at $\mu=3$ is whether $\mathbb E_w[u^{1+1/\gamma}]$ can be bounded better than by
Lyapunov, given $\mathbb E_w[u]=1$ and a bound on the variance. That is a question about moments of a
positive random variable with known mean, not about households, and it is where we would look next.

\subsection{Routes that are closed}

Three, each with a reason worth keeping so that none is reopened.

\paragraph{Mean field games.} Section~\ref{sec:mfg}. Date-by-date Lasry--Lions monotonicity is false
here at every risk aversion, and the summed condition that the argument actually uses, taken to its
weakest form, collapses to monotone capital supply. What survives is about the life-cycle household:
the reduction of the summed pairing to a mixed difference, the duration wedge for $\Phi$, and the
exact re-pricing identity.

\paragraph{Stochastic dominance.} Light and Weintraub's Theorem 2 reaches uniqueness from ordered
distributions, and the obvious weakening of its hypothesis is to ask only for that ordering: first-order
stochastic dominance of the cohort distributions. It buys nothing. In wage units, at $r=4\%$ against
$6\%$, dominance holds to machine precision at $\gamma\le3$, where the policy ordering already holds
exactly, and fails at precisely the risk aversions where the policy ordering fails --- at $\gamma=5$,
eleven of sixty ages have a crossing of the asset CDF, worst gap $5\times10^{-2}$, and $0.57$
conditional on the earnings state. The distributions genuinely cross, and they stop being ordered
exactly when the policies do, while aggregate supply stays monotone throughout. The aggregate is the
weaker statement and there is no distributional route to it (\texttt{fosdcheck.m}).

\paragraph{Two-sided sandwiches.} For the two-rate comparison at a state, the width of the sandwich at
a given rate \emph{is} the precautionary wedge: arithmetic human wealth gives the ceiling,
risk-adjusted gives the floor, and they cannot both be tight. So bounding the two problems separately
works only while the consumption response to the rate exceeds that wedge, which at Aiyagari's numbers
is $\gamma=1$ alone (Section~\ref{sec:thm1}). Any route to $\mu>1$ has to bound the \emph{difference}
between the two problems, which is what Section~\ref{sec:inc} does.

\subsection{The classical two-period results}

Section~\ref{sec:diamond} does the steady state, and for every risk aversion rather than only for
logarithmic utility. What remains of the two-period economy is standalone: monotone global
convergence, existence under Inada conditions with a multiplicity witness in general, dynamic
inefficiency of a steady state with $r<n$ by an explicit Pareto-improving transfer, the golden rule,
and the autarkic and monetary steady states of the pure-exchange economy.

\subsection{Extensions}

An age profile of earnings; a government budget alongside the pension of Section~\ref{sec:ret};
survival risk with accidental bequests as a second fixed point; population growth. Each of these
changes the equilibrium object rather than the household, so the existence argument of
Section~\ref{sec:eqm} should carry with the bracket recomputed.

\subsection{The certificate}

The uniqueness result of Section~\ref{sec:inc} rests on a hypothesis discharged by a computation in
exact rational arithmetic rather than by a Lean proof term. Two things would improve its standing. The
computation could be replayed inside Lean, which for a branch and bound over rational boxes is a
matter of \texttt{Decidable} arithmetic and a large \texttt{native\_decide}, or the variance bound could
be replaced by an analytic argument, which is the same question as the one above about moments. Until
one of those, the result has the standing of the existence bracket of Section~\ref{sec:eqm}: checked,
and not checked by the kernel.
